% Generated draft: requires mathematical and source review.
\documentclass[11pt]{article}
\usepackage[margin=25mm]{geometry}
\usepackage{fontspec}
\setmainfont{Latin Modern Roman}
\usepackage{amsmath,amssymb,mathtools,unicode-math}
\setmathfont{Latin Modern Math}
\usepackage{enumitem,needspace,xurl,hyperref}
\setlength{\parindent}{0pt}
\setlength{\parskip}{5pt}
\setlength{\emergencystretch}{3em}
\newcommand{\N}{\mathbb{N}}
\newcommand{\Z}{\mathbb{Z}}
\newcommand{\Q}{\mathbb{Q}}
\newcommand{\R}{\mathbb{R}}
\newcommand{\C}{\mathbb{C}}
\newcommand{\F}{\mathbb{F}}
\newcommand{\A}{\mathbb{A}}
\newcommand{\PP}{\mathbb{P}}
\newcommand{\E}{\mathbb{E}}
\newcommand{\eps}{\varepsilon}
\newcommand{\dd}{\,\mathrm{d}}
\begin{document}
\section*{UnsolvedMath ID: 700397}
\textbf{Source problem code:} catalogue entry
\subsection{Definitions and mathematical statement}
Let \(p\) be a prime. A \(p\)-group means a group in which every element has finite order equal to a power of \(p\); the group itself is not assumed finite. For a subgroup \(H\leq G\), write \(H^g=g^{-1}Hg\).

An internal semidirect product \(G=F\rtimes H\) means that \(F\trianglelefteq G\), \(H\leq G\), \(G=FH\), and \(F\cap H=\{1\}\). Following the supplied definition, such a group is called a Frobenius group with kernel \(F\) and complement \(H\) if
\[
H\cap H^g=\{1\}\qquad\text{for every }g\in G\setminus H
\]
and
\[
F\setminus\{1\}=G\setminus\bigcup_{g\in G}H^g.
\]

The existence question is whether there are a prime \(p\), a \(p\)-group \(G\), and subgroups \(F,H\leq G\) satisfying these conditions.

\emph{Unresolved qualification.} The supplied definition does not explicitly require \(F\neq\{1\}\) or \(H\neq\{1\}\). These nontriviality requirements affect the intended existence question and must be confirmed. If both are intended, the question is the same existence statement with the additional conditions \(F\neq\{1\}\) and \(H\neq\{1\}\). They have not been silently added to the supplied definition. The source also gives no explicit finiteness restriction; the formulation above permits infinite \(p\)-groups.
\subsection{Short English statement}
V. P. Shunkov asks whether Frobenius \(p\)-groups exist. Such a group must split into a kernel and complement satisfying the displayed conjugate-intersection and complement-union conditions. The intended nontriviality requirements for the kernel and complement need clarification.
\subsection{Sources}
\begin{itemize}
\item Local UnsolvedMath registry, permanent ID \(700397\), exact source problem code \texttt{catalogue entry}, title \emph{Kourovka Notebook Problem 6.55}. Supplied status: \texttt{open}; this is a registry label, not an independently verified mathematical assessment. Supplied upstream revision: \texttt{d371ee917ac9647fbb06cd6fc529deec6f3dda82}.
\item \emph{The Kourovka Notebook (Unsolved Problems in Group Theory)}, original problem 6.55, sixth issue (1978); proposed by V. P. Shunkov. The supplied background identifies the twenty-first issue edition (2026), \texttt{arXiv:1401.0300v48}. The supplied collection description names E. I. Khukhro and V. D. Mazurov as editors. This attribution and edition information are supplied source data and have not been independently checked.
\end{itemize}
\end{document}
